\chapter{What ``Exact'' Means}\label{ch:precision}
\index{exactness}\index{fitted model}\index{Lamb shift}\index{anomalous magnetic moment}\index{atomic units}\index{Planck's constant!as a fitted scale}

\begin{quote}
\itshape A fitted model is exact in the aspect it was fitted to.
\end{quote}
\bigskip

Every objection this book has met comes back to one comparison. Standard quantum mechanics is said
to be exact and merely uncomputable; RealQM computes but agrees with measurement only to a few per
cent. Put that way the matter is settled before it is examined. This chapter examines it, and the
sentence above is the whole of the argument. Exactness in the aspect a model was fitted to is
automatic --- it was arranged, and it carries no information. What a model does \emph{outside} its
fit is the only thing that can fail, and therefore the only thing that is evidence.

\section{The premise}\label{sec:prec-premise}

The claim that organises every introduction to the subject is that the Schr\"odinger equation
solves hydrogen exactly. For one electron about a fixed point charge,
\begin{equation}
-\tfrac12\nabla^2\psi - \frac{1}{r}\,\psi = E\psi,
\label{eq:prec-schrH}
\end{equation}
with ground state $E=-\tfrac12$ in atomic units, $13.605693$~eV, and spectrum $E_n = -1/2n^2$
reproducing the Balmer series. The agreement was decisive in 1926 and remains the reason the
equation is believed. The question is what that agreement establishes --- and to answer it one must
ask where the two terms in \eqref{eq:prec-schrH} came from, because they did not come from the same
place.

\section{Where the form came from}\label{sec:prec-form}

\paragraph{The potential is not ad hoc.}
The term $-1/r$ is Coulomb's law, and it enters \eqref{eq:prec-schrH} already established. It was
measured by Coulomb in 1785 on a torsion balance at macroscopic separations, in experiments that
knew nothing of atoms or of spectra, and the inverse-square form has since been confirmed to high
precision in settings unrelated to the hydrogen atom. It arrives with no freedom: no exponent to
select, no prefactor to tune, nothing adjusted to atomic measurements. Whatever fitting went into
\eqref{eq:prec-schrH}, the Coulomb term is not where it happened.

\paragraph{The kinetic term is.}
The operator $-\tfrac12\nabla^2$ has no comparable standing. Nothing in electrostatics or mechanics
requires the quadratic form $T[\psi] = \tfrac12\int|\nabla\psi|^2$, and it was not derived from
anything prior. Schr\"odinger reached it by analogy --- with Hamilton--Jacobi theory and with de
Broglie's relation --- and an analogy is not a derivation. The decisive detail is what happened to
the alternative: he wrote a relativistic wave equation first, found that its eigenvalues gave the
wrong fine structure for hydrogen, and set it aside in favour of the non-relativistic form, which
gave the right gross spectrum. The surviving operator is the one that produced the observed levels.
It was selected \emph{by} its spectrum.

So the two ingredients have opposite status. One carries independent empirical support obtained
without reference to atoms. The other was chosen because its eigenvalues came out as the ones
already on the table.

\paragraph{What was on the table.}
The pattern was known before the equation existed: Balmer wrote his formula for the visible
hydrogen lines in 1885 and Rydberg generalised it in 1888, both purely empirically, the lines going
as differences of $1/n^2$. Equation~\eqref{eq:prec-schrH} was accepted because it returns exactly
this. The $1/n^2$ pattern is a consequence of the equation, but the equation is a consequence of
the pattern.

The famous $-0.500000$ has a second and separate issue: its digits are not an outcome. The
coefficient $\tfrac12$, and with it a ground state at exactly $-\tfrac12$, is a statement about
units --- atomic units are \emph{defined} so that this is the number. Converting it into anything
measurable needs a bridge from atomic energy to the energy of light, and that bridge is Planck's
constant, through $\Delta E = h\nu$. But $h$ --- equivalently $R_\infty$, since the Hartree is
$2R_\infty h c$ --- is itself fixed by spectroscopy, and the hydrogen spectrum is among the
measurements that fix it. The scale is read off the spectrum and then put back into it.

In the narrow sense the equation has no adjustable parameter: given $\hbar$, $m_e$ and $e$ the
ground state follows. In the sense that matters for the word \emph{exact}, it is a model of simple
form whose form and whose scale were both selected to fit one body of data --- and it is that body
of data which is then presented as its triumph.

\section{Outside the fit}\label{sec:prec-outside}

For hydrogen the fitted aspect is the bound-state spectrum: the $1/n^2$ pattern and, through $h$,
its scale. Exactness there is free. The first genuinely unfitted predictions are close at hand, and
the equation gets them wrong.

One candidate failure must be set aside first, because it is not the equation's. The nucleus moves,
and \eqref{eq:prec-schrH} holds it fixed --- but that is a shortcoming of how the equation has been
\emph{applied}, not of the equation. Written for two bodies it separates exactly into
centre-of-mass and relative motion, and the relative part returns the hydrogenic spectrum with the
reduced mass in place of $m_e$, giving $13.598287$~eV. That is the equation's own answer, obtained
exactly. Charging the proton's motion against \eqref{eq:prec-schrH} would repeat precisely the
misattribution this chapter is about.

What remains cannot be dismissed that way. Against the measured $13.598435$~eV the residual is
$0.000147$~eV, or $1.1\times10^{-5}$ relative --- the relativistic and radiative physics the
equation does not contain. Five significant figures, not infinitely many. And the spectrum fails
more sharply, qualitatively rather than numerically: \eqref{eq:prec-schrH} makes $2s$ and $2p$
exactly degenerate, and they are not. The splitting is $1057.845$~MHz, the Lamb
shift~\cite{LambRetherford}, whose discovery in 1947 initiated renormalised quantum
electrodynamics. Here the equation does not return
a slightly wrong number; it returns zero where there is an effect.

So the equation said to be exact for hydrogen is the leading term of an expansion, completed by
Dirac and by quantum electrodynamics. None of this is news to atomic physics. It is a correction to
what the hydrogen case is taken to establish.

\section{The step that is not justified}\label{sec:prec-step}

For $N$ electrons the equation is written in $3N$ dimensions with $\Psi$ antisymmetric. This is a
separate postulate. It is not derived from \eqref{eq:prec-schrH}, and \eqref{eq:prec-schrH} does not
determine it: a one-electron equation fixes nothing about the dimension of the space the
many-electron problem is posed in. Antisymmetry is a further addition again, since the bare
equation is content with symmetric solutions and with them predicts the wrong periodic table.

The inductive support offered for the $3N$ postulate is the hydrogen case --- where, by
Section~\ref{sec:prec-form}, the exactness is in the fitted aspect and therefore empty, and where,
by Section~\ref{sec:prec-outside}, the unfitted predictions fail. That is a thin foundation for a
claim of exactness at every $N$.

\section{The showcase, and what it shows}\label{sec:prec-ae}

The claim examined here is usually made in the form Feynman gave it, and it is worth having in his
words. Having quoted the electron's magnetic moment as $1.00115965221$ from experiment against
$1.00115965246$ from theory, he writes~\cite[ch.~1]{FeynmanQED}:
\begin{quote}
\itshape At the present time I can proudly say that there is no significant difference between
experiment and theory! \dots\ To give you a feeling for the accuracy of these numbers, it comes out
something like this: If you were to measure the distance from Los Angeles to New York to this
accuracy, it would be exact to the thickness of a human hair.
\end{quote}
The figures have improved since --- the moment is now measured to $0.13$ parts per
trillion~\cite{Fan2023} and agrees with the calculated series~\cite{Aoyama2019} at about one part
in $10^{12}$ --- and this is the evidence meant whenever it is said that no observation has ever
been found to disagree with quantum mechanics. The description of it as the most precisely tested
prediction in science is fair; the inference usually drawn from it is not.

Notice first what Feynman is careful to say, and about what. His claim is about \emph{quantum
electrodynamics}, which he names in the same breath, and the number he selects is a one-particle
quantity. That is precisely the distinction this section turns on, and the canonical source for the
precision claim already observes it.

The quantity is one-particle. It concerns a single electron's coupling to the quantised
electromagnetic vacuum, computed as a perturbative series in $\alpha$ within quantum
electrodynamics. It involves no many-electron wavefunction, no antisymmetry among several
electrons, no $3N$-dimensional equation. It is overwhelming evidence that QED describes the vacuum
coupling of one charge, and it is not evidence about the many-electron equation; offering it in
support confuses two different theories that happen to share a Lagrangian. Two further points
belong with it: the comparison requires $\alpha$, whose most precise determinations now come either
from the moment itself or from atom-recoil measurements that have at times disagreed at the level
of the comparison, so part of the agreement is internal rather than independent --- and the
precision is a property of the observable, not of the theory's reach. Twelve digits on a vacuum
quantity coexists with no ability to state an error bar on a hundred-electron molecule.

The last point is not a private complaint of this book's. Leggett, reviewing what has actually been
tested, opens from the same observation~\cite[p.~R415]{Leggett2002}:
\begin{quote}
\itshape Despite this enormous range of applications, however, there is a sense in which we can say
that the region of the whole parameter space over which the validity of QM has been directly tested
is still rather modest.
\end{quote}
His concern is length scale and the size of superposed objects rather than electron count, but the
shape of the argument is the one made here: a theory can be confirmed to many digits inside a
narrow region and remain untested across most of the range over which it is nonetheless assumed.

\section{Where the equation can be solved}\label{sec:prec-solve}

For an equation advertised as exact for every $N$, the question that decides what the advertisement
is worth is: for which $N$ has it been \emph{solved}? Not approximated --- solved, to an accuracy
the equation itself controls.

The answer is two, three, and with effort four. Helium is known to some twenty digits and H$_2$ to
twelve, by explicitly correlated expansions in the interelectronic distances. These are not
truncations of a general procedure that merely becomes expensive; they are constructions built for
a handful of particles, and they do not continue.

Past that the equation is not solved. It is projected onto a finite basis and the projection is
solved instead. Full configuration interaction reaches of order ten electrons in a basis too small
to be the answer, at a cost that grows factorially, and the remaining error is the basis's rather
than the equation's --- not quantifiable against the exact solution, because the exact solution is
what is missing. Beyond ten electrons the working methods are different theories: coupled-cluster
at a chosen truncation, or density functionals carrying parameters fitted to measured data, the
latter having given up the wavefunction altogether. They are effective and they are what chemistry
uses. Neither is the Schr\"odinger equation. For a protein there is no number of any kind.

There is a further turn, and it bears on how this book's per-cent figures should be read. The
numbers a competing model is scored against are themselves largely computed. The $-2.903724$ for
helium is not a measurement: it is the solution of the nonrelativistic Schr\"odinger equation at
infinite nuclear mass, obtained by a purpose-built expansion --- the equation under test supplying
its own benchmark. What is measured is the two ionization energies, and summing the NIST values
gives $-2.903386$~Ha; the remaining $0.00034$~Ha is \emph{removed by calculation}, being the
relativistic, QED and nuclear-motion content the measurement contains and the benchmark must not.
They part company in the fourth decimal and theory supplies the difference. For atoms generally the
reference energies are explicit composites of experiment and \emph{ab initio} calculation, computed
outright past argon~\cite{Chakravorty}, because a nonrelativistic stationary-nucleus total
electronic energy is not an observable --- no instrument reads one. The composites are accurate and
this does not impugn them. It does mean that ``agrees with the exact value to $n$ per cent'' states
a distance from a construction that already contains the theory under test, which is why
Chapter~\ref{ch:atoms} rests its case on observables of another kind --- ionization energies,
geometry, shell radii, ordering.

So the equation is exact for all $N$ and solvable for $N$ of about three. An exactness testable
only where a bespoke method happens to exist, and replaced by fitted models everywhere else, has no
content at the sizes chemistry and biology are made of.

\paragraph{Why the claim of universal agreement is nearly empty.}
This is the point at which the two halves of the standard position --- \emph{exact} and
\emph{uncomputable} --- turn out to undercut one another, and it is worth stating bluntly. A
disagreement with observation requires a prediction to disagree with. Past three or four electrons
the $3N$-dimensional equation yields no number at all; the numbers come from truncations and from fitted
functionals. So for almost all of chemistry there is nothing for an experiment to contradict, and
the celebrated record of never having been contradicted is, at those sizes, not a record of success
but a record of \emph{absence of exposure}.

The practice follows accordingly, and consistently. When a computed result does fail --- dispersion
missing, a barrier wrong, a spin state out of order --- the failure is charged to the functional or
to the basis and is repaired by adjusting those, never charged to the equation. That attribution is
not dishonest; it is correct, because the equation was never evaluated and so cannot be what failed.
But it means no observation can reach the equation at all. Uncomputability places it beyond
refutation, and what is beyond refutation is equally beyond confirmation: the unbroken agreement is
a property of the arrangement, not a finding about nature.

So the claim retains content exactly where computation is possible --- hydrogen, helium, the
one-particle vacuum quantity of Section~\ref{sec:prec-ae} --- which is also where the form was fitted
and where the showcase sits. Everywhere else it is the empty statement that a theory which makes no
checkable prediction has never been caught out. This is the same structure as the principle in
Section~\ref{sec:prec-form}, arrived at from the other side: there, exactness in the fitted aspect
carried no information because it was arranged; here, agreement carries no information because
nothing was risked.


\section{What the comparison is about}\label{sec:prec-conclusion}

The two models are both postulates. Both generalise Coulomb interaction and a chosen kinetic form
to many electrons, in different ways, and neither generalisation is derived. They differ in what
can be done with them.

One is exact in a limit nobody can compute, misses by $1.1\times10^{-5}$ and by a whole splitting
in the one case it solves exactly, achieves extraordinary precision for two and three electrons by
methods that do not scale, and offers no error bar at all past ten. The other is exact nowhere,
agrees with measurement to a few per cent, and does so at every size from one electron to ten
thousand, in three dimensions, at linear cost.

The first is more accurate wherever it can be evaluated. The second is \emph{answerable} wherever
matter exists. A theory whose error is knowable and roughly constant can be held to account at
every scale; a theory whose error is unquantifiable past three electrons can be held to account
only there. Which property matters more depends on what one wants a model of matter for --- and
that is a question about scope, which is a legitimate scientific question, rather than a question
about precision, which the twelve digits have already settled in a domain where nobody disputes
them.

Chapter~\ref{ch:h2} makes the same point with numbers on the one system where both models can be
pushed to convergence: the two RealQM constructions available at $N=2$ bracket the exact energies
of He and H$_2$ from either side, and the few per cent they leave is the same few per cent the model
carries at ten thousand electrons, for the same reason.

% --- End of chapter ---
